\appendix

\section*{Artifact availability}
The benchmark, every corpus, all result files and the analysis scripts accompany this
submission as an anonymized artifact and will be released publicly on acceptance; the
recovered records of Section~\ref{sec:external} are identified by the commit hash of the
repository they were recovered from.

\section{The published comparison}
\label{app:published}

\begin{table}[h]
\centering
\begin{tabular}{lrrrr}
\toprule
Arm & Repairs & Effective & Input tokens & Unresolved boundaries \\
\midrule
\texttt{semantic\_only}  & 3 & 0 & 568{,}715   & 4 \\
\texttt{history\_full}   & 3 & 1 & 1{,}824{,}533 & 4 \\
\texttt{etiq\_full}      & 3 & 0 & 1{,}331{,}686 & 3 \\
\texttt{etiq\_selected}  & 3 & 1 & 1{,}000{,}797 & 2 \\
\bottomrule
\end{tabular}
\caption{The published comparison, one run per arm, 30 July 2026, taken from the
archived machine-readable summary. Its own record notes that ``no arm reached final
trust''.}
\label{tab:published}
\end{table}

\section{Defects in the trust machinery, in detail}
\label{app:defects}

The validator confining a repair to one boundary returned early whenever the target could
not be resolved to a function, so in that branch a repair could rewrite the whole target
file. The branch is reached whenever a declared boundary name does not resolve in the
generated source, and the authoring model writes both, so it is reachable by ordinary
error and steerable on purpose; the harness splices over the same span, so an unresolvable
target produced a whole-file substitution the validator accepted. It was untested, because
the scanner fixture reported a call stack inside functions the authoring fixture never
defined, so no target ever resolved and every workflow test of repair ran through the
unenforced branch. Separately, the module producing the published numbers was imported by
no test, resource limits around generated code defaulted to inactive, and the home
directory reached the generated pipeline, which has network access outside replay mode;
captured third-party text now reaches the reviewer inside a digest-keyed fence it cannot
close \cite{greshake2023injection}.

\section{A declared invariant, concretely}
\label{app:invariant}

The retrieval stage of the branching pipeline carries: \texttt{retrieve\_registries} must
return \texttt{source\_ids} and \texttt{source\_class}; \texttt{source\_ids} must hold at
least three entries; every entry must match \texttt{\^{}[a-z]+\_[a-z]+\$}. Three clauses
catching a disappeared field, a silently shortened retrieval, and a malformed identifier;
what none can catch is an identifier that is correctly shaped and simply wrong, which is
measured rather than assumed in the body.

\section{Extended related work}
\label{app:related}

Empirical work on why multi-agent systems fail supplies the closest framing for our
problem: Cemri et al.\ build a taxonomy from more than two hundred annotated traces and
make task verification one of its three top-level categories \cite{cemri2025multiagent},
though their failures are annotated by human experts rather than detected by the system at
runtime; Chen et al.\ survey the same territory as a gap between what long-horizon agents
are asked to do and what their planning, memory and evaluation machinery supports
\cite{chen2026horizon}. The assumption our benchmark tests is that a language model
supplies that verification, and there is reason for caution: judges favour text more
familiar to themselves \cite{wataoka2024selfpreference}, which matters because reviewer
and reviewed here share a model family; a judge's self-agreement can be high while its
agreement with the truth is not \cite{norman2026reliability}, the distinction our
localisation metric operationalises; verifiers checking only what they can express admit
confident false positives \cite{helff2026gaming}; and multi-step agents exploit
naturalistic chances to tamper with what evaluates them \cite{thaman2026rewardhacking}.
Provenance-based trust for agents \cite{wang2026traces} and recoverable execution from
recorded checkpoints \cite{zhuang2026agentrewind} are the constructive lines this work
measures rather than proposes.

Our method borrows its central move from mutation testing, a defect of known location
measuring whether a suite detects it, including the equivalent-mutant problem we meet as
an injection that changes nothing \cite{jia2011mutation}. The substrate is execution
provenance \cite{freire2008provenance}, whose distinction between a record of what a
computation did and a claim about what it meant locates this paper's difficulty: the graph
is a record, the trust label is a claim, and only the first is checkable. The fencing of
Section~\ref{sec:defects} addresses indirect prompt injection through fetched content
\cite{greshake2023injection}, a designed-in exposure here since the pipeline must retrieve
third-party sources at runtime.

\section{Staleness intervention, full results}
\label{app:stale}


We crossed two factors, twenty cases per cell on each of two judges: the history shown
(clean prior iterations, or a superseded run carrying a fault since repaired) and the
current run (clean, or carrying a new fault at a different boundary), asking for the
faulty stage of the current run or none.

\begin{table}[h]
\centering
\begin{tabular}{llrr}
\toprule
History & Current run & Judge A & Judge B \\
\midrule
clean & clean     & 18/20 & 4/20 \\
stale & clean     & 17/20 & 13/20 \\
clean & new fault & 18/20 & 19/20 \\
stale & new fault & 11/20 & 19/20 \\
\bottomrule
\end{tabular}
\caption{Correct decisions under the staleness intervention. Each judge fails in a
different place: A dismisses genuine new faults when a stale prior is present, B cannot
recognise a clean run unless one is.}
\label{tab:stale}
\end{table}

Neither judge re-flagged the repaired boundary, which was the predicted failure; that
happened once in 160 stale-history calls. Judge A, handed a superseded faulty prior,
misses genuine new faults, 18/20 falling to 11/20 (Fisher exact, $p=0.03$ uncorrected,
suggestive after the multiplicity correction of Section~\ref{sec:power}); the prior appears
to act as a contrast anchor, and a current run that looks better than the old one reads as
fine. Judge B shows the opposite defect: it detects the new fault regardless of history,
19/20 in both cells, but cannot recognise a clean run, inventing a fault in 16 of 20
clean-history cases and naming the final stage in every one, which is the downstream-blame
bias with nothing to blame; a superseded faulty prior partially restores its ability to
answer clean, 13/20 against 4/20. The deployed system's own record supplies a third
signature, a frontier model resurrecting a resolved issue from flat history.

Three judges, three failure shapes from one manipulation. The stable conclusion is a
dependence rather than a direction: the content of accumulated history modulates a model
judge's trust decisions in model-specific ways that no prompt here controlled, while
deterministic signals are unaffected by context by construction.
